% European Heart Journal Top 10 VHD 2025 教學講義
% 2025年瓣膜疾病年度十大論文回顧
\documentclass[12pt,a4paper]{article}
% Drake_preamble.tex - LaTeX Preamble Template
% 作者: 謝慕揚, MD, PhD, FESC
% 
% 使用說明:
% 1. 表格請使用 longtable，不要有垂直分隔線 |
% 2. 表格內換行使用 \newline，不要使用 \\
% 3. 藥物名稱、檢驗、檢查請維持英文
% 4. Reference 格式請使用 \href 加上驗證過的 DOI

% Including essential packages for Chinese typesetting
\usepackage{xeCJK}
\usepackage{fontspec}
% Configuring Chinese font with Noto Serif CJK TC, fallback to SimSun
\IfFontExistsTF{Noto Serif CJK TC}{
  \setCJKmainfont{Noto Serif CJK TC}
  \setCJKsansfont{Noto Serif CJK TC}
  \setCJKmonofont{Noto Serif CJK TC}
}{\setCJKmainfont{SimSun}
  \setCJKsansfont{SimSun}
  \setCJKmonofont{SimSun}
}

% Including other necessary packages
\usepackage{geometry}
\usepackage{titlesec}
\usepackage{enumitem}
\usepackage{xcolor}
\usepackage{colortbl}
\usepackage{tcolorbox}
\usepackage{booktabs}
\usepackage{longtable}
\usepackage{array}
\usepackage{graphicx}
\usepackage{tikz}
\usepackage{fancyhdr}
\usepackage{multicol}
\usepackage{datetime2}
\usepackage{hyperref}
\usepackage{mdframed}
\usepackage{amsmath}
\usepackage{amssymb}
\usepackage{multirow}

\hypersetup{
    colorlinks=true,
    linkcolor=emergency_blue,
    citecolor=emergency_blue,
    urlcolor=emergency_blue,
    pdfborder={0 0 0}
}

% Defining page geometry
\geometry{
  top=2.5cm,
  bottom=2.5cm,
  left=2cm,
  right=2cm,
  headheight=25pt
}

% Adjusting topmargin to compensate for increased headheight
\addtolength{\topmargin}{-10pt}

% Defining colors
\definecolor{emergency_red}{RGB}{186,24,27}
\definecolor{emergency_blue}{RGB}{0,114,188}
\definecolor{emergency_gray}{RGB}{120,120,120}
\definecolor{highlight_yellow}{RGB}{255,250,205}

% Setting title formats
\titleformat{\section}[block]{\Large\bfseries\color{emergency_red}}{\thesection}{10pt}{}
\titleformat{\subsection}[block]{\large\bfseries\color{emergency_blue}}{\thesubsection}{8pt}{}
\titleformat{\subsubsection}[block]{\normalsize\bfseries\color{emergency_gray}}{\thesubsubsection}{6pt}{}

% Defining custom environment for important notes
\newmdenv[
    linecolor=emergency_red,
    backgroundcolor=highlight_yellow,
    linewidth=2pt,
    topline=true,
    bottomline=true,
    leftline=true,
    rightline=true,
    innertopmargin=10pt,
    innerbottommargin=10pt,
    innerleftmargin=10pt,
    innerrightmargin=10pt
]{importantbox}

% Defining note command with tcolorbox
\newcommand{\note}[1]{
    \par
    \noindent
    \begin{tcolorbox}[
        colback=emergency_blue!5!white,
        colframe=emergency_blue,
        title=\textbf{重點提醒},
        fonttitle=\bfseries,
        boxrule=0.5pt,
        arc=3pt,
        left=6pt,
        right=6pt,
        top=6pt,
        bottom=6pt
    ]
    #1
    \end{tcolorbox}
}

% Key findings box
\newcommand{\keyfinding}[1]{
    \par
    \noindent
    \begin{tcolorbox}[
        colback=yellow!10!white,
        colframe=orange!75!black,
        title=\textbf{核心發現摘要},
        fonttitle=\bfseries,
        boxrule=0.5pt,
        arc=3pt
    ]
    #1
    \end{tcolorbox}
}

% Defining custom date format for ISO 8601
\DTMnewdatestyle{iso}{
  \renewcommand{\DTMdisplaydate}[4]{\number##1-\number##2-\number##3}
  \renewcommand{\DTMDisplaydate}[4]{\number##1-\number##2-\number##3}
}
\DTMsetdatestyle{iso}
\newcommand{\mydate}{\DTMtoday}


% Custom header for this document
\pagestyle{fancy}
\fancyhf{}
\fancyhead[L]{\textcolor{emergency_red}{European Heart Journal 教學講義}}
\fancyhead[R]{\textcolor{emergency_blue}{瓣膜疾病年度回顧}}
\fancyfoot[C]{\textcolor{emergency_gray}{\thepage}}
\renewcommand{\headrulewidth}{1pt}
\renewcommand{\headrule}{\hbox to\headwidth{\color{emergency_red}\leaders\hrule height \headrulewidth\hfill}}

\begin{document}

%============================================================
% Title Page
%============================================================
\begin{center}
{\LARGE\bfseries\textcolor{emergency_red}{European Heart Journal}}\\[0.5cm]
{\Large\textbf{The Year in Cardiovascular Medicine 2025:\\The Top 10 Papers in Valvular Heart Disease}}\\[0.3cm]
{\large 2025年心血管醫學年度回顧：瓣膜性心臟病十大論文}\\[0.5cm]
{\normalsize \href{https://doi.org/10.1093/eurheartj/ehaf1088}{Eur Heart J 2026. DOI: 10.1093/eurheartj/ehaf1088}}\\[0.3cm]
{\normalsize 整理：謝慕揚 MD, PhD, FESC \quad 日期：\mydate}
\end{center}

\vspace{0.5cm}

\keyfinding{
本文由三位專家（Baumgartner、Iung、Messika-Zeitoun）從 NEJM、JAMA、Lancet、Nature、EHJ、Circulation、JACC 中精選 2025 年瓣膜疾病領域最具影響力的十篇論文，涵蓋：
\begin{itemize}[noitemsep]
    \item 人工智慧（AI）在瓣膜疾病診斷與風險分層的應用
    \item 無症狀嚴重主動脈瓣狹窄（AS）的處置爭議
    \item TAVI vs 手術在低風險患者的長期追蹤
    \item TAVI 後藥物治療（Dapagliflozin）
    \item 二尖瓣與三尖瓣介入治療的最新證據
\end{itemize}
}

%============================================================
\section{人工智慧在瓣膜疾病的應用}
%============================================================

\subsection{ECHONEXT 研究：AI 偵測結構性心臟病}

\begin{itemize}
    \item \textbf{研究設計}：Deep learning ECG model，訓練資料超過一百萬筆心律與影像紀錄
    \item \textbf{目標}：偵測中重度瓣膜疾病（VHD）等結構性心臟病
    \item \textbf{結果}：
    \begin{itemize}
        \item 內部與外部驗證均展現高診斷準確度
        \item 在不同照護場域及種族群體表現一致
        \item 臨床試驗中成功識別先前未診斷的心臟病患者
    \end{itemize}
    \item \textbf{意義}：有潛力擴展大規模心臟病篩檢的可近性
\end{itemize}

\subsection{DELINEATE-REGURGITATION 研究：AI 評估瓣膜逆流}

\begin{itemize}
    \item \textbf{資料}：71,660 份經胸心臟超音波，1,203,980 段 color Doppler 影像
    \item \textbf{準確度}（Weighted kappa，內部/外部驗證）：
    \begin{itemize}
        \item Aortic regurgitation (AR)：0.81 / 0.76
        \item Mitral regurgitation (MR)：0.76 / 0.72
        \item Tricuspid regurgitation (TR)：0.73 / 0.64
    \end{itemize}
    \item \textbf{特色}：多視角方法優於單視角；DELINEATE-MR-Progression 可預測 MR 進展
\end{itemize}

\note{這些研究顯示 AI 即將在瓣膜疾病的臨床管理中產生重大影響。}

%============================================================
\section{無症狀嚴重主動脈瓣狹窄的處置}
%============================================================

\subsection{四項 RCT 的 Meta-analysis}

\begin{longtable}{p{4cm}p{10cm}}
\toprule
\textbf{項目} & \textbf{內容} \\
\midrule
\endfirsthead
\midrule
\textbf{項目} & \textbf{內容} \\
\midrule
\endhead
\bottomrule
\endfoot
納入研究 & 4 項 RCT，共 1,427 位患者 \newline 比較早期 AVR/TAVI vs 臨床監測 \\
\midrule
納入條件 & 無症狀嚴重高壓差 AS \newline 左心室功能保留 \\
\midrule
平均追蹤 & 4.1 年 \\
\midrule
主要發現 & 早期 AVR 顯著降低：\newline • 非計畫性 CV/HF 住院（14.6\% vs 31.9\%, P < .01）\newline • 中風（4.5\% vs 7.2\%, P = .03）\\
\midrule
死亡率 & 全因及心血管死亡：無顯著差異 \newline 但研究間存在高度異質性 \\
\midrule
爭議 & 最大的 RCT 中，監測組於 6 個月內交叉至介入被計為事件 \newline 此定義頗具爭議 \\
\midrule
需考量 & 需權衡短中期獲益與長期風險（人工瓣膜相關併發症、再介入） \\
\bottomrule
\end{longtable}

\note{DANAVR 和 EASY-AS 等進行中的隨機試驗將提供更多證據。}

%============================================================
\section{低風險患者：TAVI vs 外科手術}
%============================================================

\subsection{PARTNER-3 試驗：7 年追蹤}

\begin{itemize}
    \item \textbf{設計}：1,000 位低風險手術患者隨機分配至 TAVI 或 SAVR
    \item \textbf{7 年結果}：
    \begin{itemize}
        \item 兩個主要複合終點（死亡、中風、再住院）無顯著差異
        \item 血流動力學結果與瓣膜衰竭率相當
        \item 值得注意：全因死亡的 Kaplan-Meier 曲線在 3 年後交叉
        \item TAVI 組數值上較高死亡率（19.5\% vs 16.8\%）
    \end{itemize}
\end{itemize}

\subsection{NOTION-2 試驗：3 年追蹤}

\begin{itemize}
    \item \textbf{特色}：370 位低風險患者，平均年齡 71 歲（迄今最年輕）
    \item \textbf{首次納入}：雙葉主動脈瓣（Bicuspid aortic valve, BAV）
    \item \textbf{3 年結果}：
    \begin{itemize}
        \item 整體主要終點：TAVI 16.1\% vs SAVR 12.6\%（P = .4）
        \item \textbf{三葉瓣 AS}：兩組相近（14.5\% vs 14.4\%, P = 1.0）
        \item \textbf{雙葉瓣 AS}：TAVI 風險較高（20.4\% vs 7.8\%, P = .08）
    \end{itemize}
\end{itemize}

\begin{importantbox}
\textbf{臨床建議：}
\begin{itemize}[noitemsep]
    \item 三葉瓣低風險患者：TAVI 與手術結果相當
    \item \textbf{雙葉瓣患者}：應謹慎，需更多研究
    \item 年輕患者需考量終生管理：預期再介入、TAVI 取出等
\end{itemize}
\end{importantbox}

%============================================================
\section{TAVI 後的藥物治療：DapaTAVI 試驗}
%============================================================

\begin{longtable}{p{4cm}p{10cm}}
\toprule
\textbf{項目} & \textbf{內容} \\
\midrule
\endfirsthead
\midrule
\textbf{項目} & \textbf{內容} \\
\midrule
\endhead
\bottomrule
\endfoot
設計 & 開放標籤 RCT \newline TAVI 後隨機分配至 Dapagliflozin vs 標準照護 \\
\midrule
患者 & 620 vs 637 位 \newline 平均年齡 82 歲，49\% 女性 \\
\midrule
納入條件 & 心衰竭病史 + 至少一項：\newline • 腎功能不全\newline • 糖尿病\newline • LVEF $\leq$40\% \\
\midrule
主要終點 & 1 年全因死亡或心衰竭惡化（住院或緊急就醫）\\
\midrule
結果 & Dapagliflozin 15.0\% vs 標準照護 20.1\%（P = .02）\\
\midrule
驅動因素 & 主要由心衰竭惡化降低驅動 \newline 死亡率無差異（7.8\% vs 8.9\%）\\
\bottomrule
\end{longtable}

\note{此研究強調即使在瓣膜介入後，最佳化藥物治療仍然重要。}

%============================================================
\section{二尖瓣逆流：手術與介入治療}
%============================================================

\subsection{MITRACURE Registry：二尖瓣手術真實世界數據}

\begin{itemize}
    \item \textbf{族群}：連續接受單獨或合併二尖瓣手術的 MR 患者
    \item \textbf{患者特徵}：
    \begin{itemize}
        \item 最常見病因：Myxomatous（61\%）
        \item NYHA Class III/IV：43\%
        \item 心衰竭：30\%
        \item 心房顫動/flutter：22\%
        \item EF 下降：35\%
        \item 肺高壓：22\%
    \end{itemize}
    \item \textbf{早期介入}：僅 3\% 接受早期介入
    \item \textbf{修補率}：整體 62\%，退化性 MR 80\%
    \item \textbf{院內死亡率}：整體 4.5\%，退化性 MR 2.3\%（單獨手術 1.8\%）
\end{itemize}

\note{此研究凸顯患者仍常被延遲轉介，早期介入比例低，需改善及時處置策略。}

\subsection{OCEAN-Mitral Registry：心房功能性 MR 的 TEER}

\begin{itemize}
    \item \textbf{背景}：心房功能性 MR（Atrial functional MR）被排除於三項現有 RCT 之外
    \item \textbf{比較}：TEER（441 人）vs 藥物治療（640 人，來自 REVEAL-AFMR registry）
    \item \textbf{結果}：
    \begin{itemize}
        \item TEER 降低全因死亡 + HF 住院複合終點（HR 0.65, P = .044）
        \item TEER 降低全因死亡（HR 0.58, P = .044）
        \item \textbf{關鍵發現}：僅殘餘 MR $\leq$ mild 者獲益
        \item 殘餘 MR $\geq$ moderate 者結果與藥物治療相近
    \end{itemize}
\end{itemize}

\subsection{SAPIEN M3 Trial：經導管二尖瓣置換}

\begin{itemize}
    \item \textbf{適應症}：不適合手術或 TEER 的嚴重 MR 患者
    \item \textbf{篩選}：1,171 人篩選，299 人植入（1/3 primary MR，2/3 secondary MR）
    \item \textbf{30 天結果}：
    \begin{itemize}
        \item 死亡率 0.7\%（STS score 6.6\%）
        \item 無術中死亡
        \item 無 LVOT obstruction 導致血流動力學影響
        \item 無轉為手術
    \end{itemize}
    \item \textbf{主要終點}：全因死亡 + HF 住院 25.2\%，顯著低於預設目標 45\%
    \item \textbf{限制}：篩選失敗率高，無對照組
\end{itemize}

%============================================================
\section{三尖瓣逆流：TRILUMINATE Pivotal 2 年結果}
%============================================================

\begin{longtable}{p{4cm}p{10cm}}
\toprule
\textbf{項目} & \textbf{內容} \\
\midrule
\endfirsthead
\midrule
\textbf{項目} & \textbf{內容} \\
\midrule
\endhead
\bottomrule
\endfoot
設計 & 572 位嚴重續發性 TR 患者 \newline TEER + 藥物 vs 單純藥物 \\
\midrule
1 年結果 & 生活品質：TEER 較佳 \newline 死亡或 HF 住院：無差異 \\
\midrule
2 年結果（允許交叉） & HF 住院年化率：TEER 0.19 vs 藥物 0.26（P = .02）\newline 全因死亡：無差異（17.9\% vs 17.1\%）\newline 三尖瓣手術率：無差異（2.3\% vs 4.3\%）\\
\midrule
限制 & 無 sham 對照的開放試驗 \newline 未提供 time-to-event 分析 \\
\bottomrule
\end{longtable}

\note{對於嚴重續發性 TR，仍需進一步優化 TEER 的患者選擇標準。}

%============================================================
\section{2025 年瓣膜疾病研究摘要圖}
%============================================================

\begin{center}
\begin{tikzpicture}[
    node distance=0.8cm,
    box/.style={rectangle, draw=emergency_blue, fill=emergency_blue!10,
                text width=4.5cm, minimum height=1.2cm, align=center,
                rounded corners, font=\small}
]
    \node[box, fill=emergency_red!20, draw=emergency_red] (ai) {AI 診斷\\ECHONEXT\\DELINEATE-REGURGITATION};
    \node[box, right=0.5cm of ai] (as) {主動脈瓣狹窄\\Meta-analysis\\PARTNER-3 / NOTION-2\\DapaTAVI};
    \node[box, right=0.5cm of as] (mr) {二尖瓣逆流\\MITRACURE\\OCEAN-Mitral\\SAPIEN M3};
    \node[box, below=0.5cm of as] (tr) {三尖瓣逆流\\TRILUMINATE 2-year};
\end{tikzpicture}
\end{center}


%============================================================
\section{臨床要點總結}
%============================================================

\begin{enumerate}
    \item \textbf{AI}：ECG-based deep learning 可偵測結構性心臟病，AI 評估瓣膜逆流準確度高
    \item \textbf{無症狀嚴重 AS}：早期介入可降低住院和中風，但死亡率無差異；需等待更多長期數據
    \item \textbf{低風險 TAVI}：三葉瓣患者 TAVI vs 手術結果相當；\textbf{雙葉瓣應謹慎}
    \item \textbf{TAVI 後藥物}：Dapagliflozin 可降低高風險患者的心衰竭惡化
    \item \textbf{MR 手術}：患者常延遲轉介，早期介入比例僅 3\%
    \item \textbf{心房功能性 MR}：TEER 可能有益，但需達到良好修復（殘餘 MR $\leq$ mild）
    \item \textbf{經導管二尖瓣置換}：對不適合手術/TEER 者是可行選項
    \item \textbf{TR TEER}：可降低 HF 住院，但對死亡率無影響
\end{enumerate}

%============================================================
\section{參考文獻}
%============================================================

\begin{enumerate}
    \item Baumgartner H, Iung B, Messika-Zeitoun D. The year in cardiovascular medicine 2025: the top 10 papers in valvular heart disease. \href{https://doi.org/10.1093/eurheartj/ehaf1088}{\textit{Eur Heart J}. 2026. DOI: 10.1093/eurheartj/ehaf1088}

    \item Poterucha TJ, et al. Detecting structural heart disease from electrocardiograms using AI. \href{https://doi.org/10.1038/s41586-025-09227-0}{\textit{Nature}. 2025;644:221-30.}

    \item Long A, et al. Deep learning for echocardiographic assessment and risk stratification of aortic, mitral, and tricuspid regurgitation: the DELINEATE-regurgitation study. \href{https://doi.org/10.1093/eurheartj/ehaf248}{\textit{Eur Heart J}. 2025;46:2780-91.}

    \item Généreux P, et al. Aortic valve replacement vs clinical surveillance in asymptomatic severe aortic stenosis: a systematic review and meta-analysis. \href{https://doi.org/10.1016/j.jacc.2024.11.006}{\textit{J Am Coll Cardiol}. 2025;85:912-22.}

    \item Leon MB, et al. Transcatheter or surgical aortic-valve replacement in low-risk patients at 7 years. \href{https://doi.org/10.1056/NEJMoa2509766}{\textit{N Engl J Med}. 2025.}

    \item Jørgensen TH, et al. Three-year follow-up of the NOTION-2 trial: TAVR versus SAVR to treat younger low-risk patients with tricuspid or bicuspid aortic stenosis. \href{https://doi.org/10.1161/CIRCULATIONAHA.125.076678}{\textit{Circulation}. 2025;152:1326-37.}

    \item Raposeiras-Roubin S, et al. Dapagliflozin in patients undergoing transcatheter aortic-valve implantation. \href{https://doi.org/10.1056/NEJMoa2500366}{\textit{N Engl J Med}. 2025;392:1396-405.}

    \item Messika-Zeitoun D, et al. Clinical presentation and outcomes after surgery for mitral regurgitation: real-world insights from the MITRACURE international registry. \href{https://doi.org/10.1161/CIRCULATIONAHA.124.073674}{\textit{Circulation}. 2025;152:927-38.}

    \item Kaneko T, et al. Transcatheter edge-to-edge repair vs medical therapy in atrial functional mitral regurgitation. \href{https://doi.org/10.1093/eurheartj/ehaf511}{\textit{Eur Heart J}. 2025.}

    \item Guerrero ME, et al. Percutaneous transcatheter valve replacement in individuals with mitral regurgitation unsuitable for surgery or transcatheter edge-to-edge repair. \href{https://doi.org/10.1016/S0140-6736(25)02073-2}{\textit{Lancet}. 2025;406:2541-50.}

    \item Kar S, et al. Two-year outcomes of transcatheter edge-to-edge repair for severe tricuspid regurgitation: the TRILUMINATE pivotal randomized controlled trial. \href{https://doi.org/10.1161/CIRCULATIONAHA.125.074536}{\textit{Circulation}. 2025;151:1630-8.}
\end{enumerate}

\end{document}
